\chapter{Metodología}\label{cap:metodologia}

En este capítulo se describe el marco metodológico de este trabajo de tesis.

\section{El conjunto de datos MessIRve}\label{sec:metodologia-dataset}
% - Cifras: consultas, documentos, juicios, variedades dialectales
% - Distribución de relevantes por consulta (pendiente en notes.md; justifica la batería de métricas)
% - Partición usada (entrenamiento/prueba por artículo de Wikipedia; restricciones del conjunto de prueba)

\section{Sistema de recuperación}\label{sec:metodologia-sistema}
% - Arquitectura del pipeline: primera etapa + fusión + reordenamiento (figura)
% - Configuración de cada modelo base; nota sobre el límite de 0.6B
% - Rerankers y profundidad de reordenamiento
% - Implementación y reproducibilidad (ir-spanish, hardware, semillas)

\section{Algoritmos de fusión}\label{sec:metodologia-fusion}
% - Los seis con sus parámetros: CombMNZ, BordaFuse, Condorcet, RRF, ISR (Inverse Square Rank), RBC (Rank-Biased Centroids)
% - Qué listas se fusionan y profundidad de corte
% - Normalización de puntuaciones para CombMNZ

\section{Diseño experimental}\label{sec:metodologia-diseno}
% - Protocolo común: métricas (nDCG@10 primaria; Recall@100, MRR@10, MAP, P@10/50) y pruebas de significancia

\subsection{Experimento 1: fusión vs. modelos individuales}\label{subsec:metodologia-e1}
% - Cada fusión vs. cada modelo individual

\subsection{Experimento 2: comparación de algoritmos de fusión}\label{subsec:metodologia-e2}
% - Comparación cruzada de los seis algoritmos, por métrica y por familia

\subsection{Experimento 3: fusión vs. reordenamiento neuronal}\label{subsec:metodologia-e3}
% - Mejor fusión vs. ambos rerankers, con análisis de eficiencia

\subsection{Experimento 4: modelos abiertos vs. propietarios}\label{subsec:metodologia-e4}
% - Sistemas de la tesis vs. líneas base propietarias reportadas en MessIRve
